\topic{5}{Why a momentum solution needs a pressure correction}
Suppose the velocity and pressure estimates have already been updated $k$
times. Superscript $k$ labels that steady coupling iteration. It is not a
physical time level. The current pressure column is $p^k$.

Hold the momentum coefficients fixed and solve
\[
Aq^*=b-Gp^k.
\]
The star labels the \emph{predictor}: the velocity obtained using this pressure
estimate. It is provisional because the momentum solve alone does not impose
continuity in every cell. Compute its cell residuals from $Dq^*$.

For one square pressure cell with $h=0.25$, suppose its predicted velocities are
\[
(u_e^*,u_w^*,v_n^*,v_s^*)=(0.12,0.08,0.01,0.03).
\]
The star on each component has the same predictor meaning. The integrated
outward volume flow rate per unit depth, called $Q_P^*$, is
\[
Q_P^*=h(u_e^*-u_w^*+v_n^*-v_s^*)
=0.25(0.12-0.08+0.01-0.03)=0.005.
\]
Dividing by cell area gives the continuity residual,
\[
r_{c,P}^*=Q_P^*/h^2=0.005/0.0625=0.08.
\]
More fluid leaves than enters. The face velocities must change so that their
combined outward flow changes by $-0.005$. Pressure correction is obtained
from this conservation requirement over all cells, rather than guessed from
this number alone.

Let $p_{\rm corr}$ denote the pressure sought by this correction solve, and
$q_{\rm corr}$ its corrected velocity. Define their increments by
\[
p'=p_{\rm corr}-p^k,\qquad q'=q_{\rm corr}-q^*.
\]
A prime denotes an additive change, not a derivative. Equivalently,
\[
\boxed{\Delta p=p_{\rm corr}-p^k=p'.}
\]
This is a pressure change at a fixed cell between estimates.
The difference $p_E-p_P$ compares pressures at two locations in one estimate.
The residual $r_{c,P}^*$ instead measures a velocity imbalance. These three
quantities have different roles and must not be interchanged.
\newpage
\topic{5 continued}{Derive the velocity correction from one momentum row}
For horizontal face $e$, the predictor equation is
\[
a_eu_e^*=\sum_{g\in\mathrm{NB}(e)}a_{eg}q_g^*
          +b_e-\frac{p_E^k-p_P^k}{h}.
\]
Use the same coefficients for a corrected solution. Substitute
$u_{e,\rm corr}=u_e^*+u'_e$, $q_{g,\rm corr}=q_g^*+q'_g$ and
$p_{P,\rm corr}=p_P^k+p'_P$, with the analogous expression at $E$:
\begin{align*}
a_e(u_e^*+u'_e)={}&\sum_{g\in\mathrm{NB}(e)}a_{eg}(q_g^*+q'_g)+b_e\\
 &-\frac{(p_E^k+p'_E)-(p_P^k+p'_P)}{h}.
\end{align*}
Subtract the predictor row. The known source $b_e$, predictor neighbor terms
and old pressure gradient cancel, leaving
\[
a_eu'_e=\sum_{g\in\mathrm{NB}(e)}a_{eg}q'_g-\frac{p'_E-p'_P}{h}.
\]
This exact correction equation couples corrections at neighboring velocity
locations. For all velocity rows it is $Aq'=-Gp'$.

\textbf{The SIMPLE approximation.} For the local velocity correction,
neglect the neighbor-correction sum. The center response remains:
\[
a_eu'_e\simeq-\frac{p'_E-p'_P}{h},\qquad
\boxed{u'_e=-d_e\frac{p'_E-p'_P}{h}},\qquad d_e=1/a_e.
\]
This approximates the momentum response to a pressure increment. It does not
remove the requirement to conserve mass. The pressure equation is chosen so
that the resulting face corrections satisfy every cell's continuity balance.

The same construction at the other three faces gives
\begin{align*}
u'_w&=-d_w\frac{p'_P-p'_W}{h},\\
v'_n&=-d_n\frac{p'_N-p'_P}{h},&
v'_s&=-d_s\frac{p'_P-p'_S}{h}.
\end{align*}
Each $d$ is the reciprocal diagonal of that face's own momentum row. Corrected
continuity on a square cell requires
\[
0=r_{c,P}^*+\frac{u'_e-u'_w+v'_n-v'_s}{h}.
\]
Section 6 substitutes the four expressions into this balance to obtain
the pressure equation. Repeated coupling iterations also restore consistency
with the full momentum equations, whose neighbor corrections were approximated.
